\documentclass[11pt,a4paper]{article}

% ============================================================
% CORRIGE DS ANALYSE PCSI -- SUITES, CONTINUITE, DERIVATION
% Niveau solide PCSI / début CCP-INP
% ============================================================

\usepackage[utf8]{inputenc}
\usepackage[T1]{fontenc}
\usepackage[french]{babel}
\usepackage{lmodern}
\usepackage{microtype}
\usepackage{geometry}
\geometry{margin=1.7cm}

\usepackage{amsmath,amssymb,amsthm}
\usepackage{mathtools}
\usepackage{array}
\usepackage{tabularx}
\usepackage{booktabs}
\usepackage{enumitem}
\usepackage{xcolor}
\usepackage{fancyhdr}
\usepackage{lastpage}
\usepackage{titlesec}
\usepackage{tcolorbox}
\usepackage{hyperref}

\hypersetup{
    colorlinks=true,
    linkcolor=blue!60!black,
    urlcolor=blue!60!black
}

% ============================================================
% COULEURS
% ============================================================

\definecolor{mainblue}{RGB}{30,70,130}
\definecolor{darkred}{RGB}{140,30,30}
\definecolor{darkgreen}{RGB}{20,110,70}
\definecolor{lightblue}{RGB}{235,242,252}
\definecolor{lightred}{RGB}{252,238,238}
\definecolor{lightgreen}{RGB}{238,250,244}
\definecolor{lightgray}{RGB}{245,245,245}

% ============================================================
% EN-TETES
% ============================================================

\pagestyle{fancy}
\fancyhf{}
\lhead{\textsc{PCSI -- Analyse}}
\chead{\textsc{Corrigé DS : Suites, continuité, dérivation}}
\rhead{Page \thepage/\pageref{LastPage}}
\cfoot{}

\titleformat{\section}
{\Large\bfseries\color{mainblue}}
{\thesection}{0.8em}{}

\titleformat{\subsection}
{\large\bfseries\color{darkred}}
{\thesubsection}{0.8em}{}

\titleformat{\subsubsection}
{\normalsize\bfseries\color{darkgreen}}
{\thesubsubsection}{0.8em}{}

% ============================================================
% ENVIRONNEMENTS
% ============================================================

\newtcolorbox{methodebox}{
    colback=lightblue,
    colframe=mainblue,
    title=\textbf{Méthode},
    fonttitle=\bfseries,
    arc=2mm,
    boxrule=0.8pt
}

\newtcolorbox{remarquebox}{
    colback=lightgreen,
    colframe=darkgreen,
    title=\textbf{Remarque},
    fonttitle=\bfseries,
    arc=2mm,
    boxrule=0.8pt
}

\newtcolorbox{attentionbox}{
    colback=lightred,
    colframe=darkred,
    title=\textbf{Point de vigilance},
    fonttitle=\bfseries,
    arc=2mm,
    boxrule=0.8pt
}

\newtcolorbox{resultatbox}{
    colback=lightgray,
    colframe=black!65,
    title=\textbf{Résultat obtenu},
    fonttitle=\bfseries,
    arc=2mm,
    boxrule=0.8pt
}

\newcommand{\R}{\mathbb{R}}
\newcommand{\N}{\mathbb{N}}
\newcommand{\eps}{\varepsilon}

\setlist[itemize]{label=\(\triangleright\), leftmargin=1.2cm}
\setlist[enumerate]{leftmargin=1cm}

% ============================================================
% DOCUMENT
% ============================================================

\begin{document}

\begin{center}
    {\Huge \bfseries Corrigé détaillé}\\[0.3cm]
    {\Large \bfseries Devoir surveillé d'analyse}\\[0.2cm]
    {\large Suites, continuité, dérivation}\\[0.4cm]
    \rule{0.85\linewidth}{0.5pt}\\[0.2cm]
    {\large Niveau PCSI -- esprit CCP-INP}\\[0.2cm]
    \rule{0.85\linewidth}{0.5pt}
\end{center}

\vspace{0.4cm}



\newpage

% ============================================================
% PROBLEME I
% ============================================================

\section*{Problème I -- La racine de \(2\) comme objet d'ordre, puis comme limite d'un algorithme}

On pose
\[
A=\{x\in\R_+ \mid x^2\leq 2\}.
\]

Le but est d'abord de construire rigoureusement \(\sqrt2\) comme borne supérieure d'un ensemble, puis d'étudier la suite de Héron.

% ------------------------------------------------------------

\subsection*{Partie A -- Construction par borne supérieure}

\subsubsection*{1. Montrer que \(A\) est non vide et majoré}

On remarque que
\[
0\in\R_+
\qquad\text{et}\qquad
0^2=0\leq 2.
\]
Donc
\[
0\in A.
\]
Ainsi, \(A\) est non vide.

Montrons maintenant que \(A\) est majoré. Soit \(x\in A\). Alors
\[
x\geq 0
\qquad\text{et}\qquad
x^2\leq 2.
\]
Si \(x>2\), alors, comme \(x\geq 0\),
\[
x^2>4>2,
\]
ce qui contredit \(x^2\leq 2\). Donc nécessairement
\[
x\leq 2.
\]
Ainsi,
\[
\forall x\in A,\qquad x\leq 2.
\]
Donc \(2\) est un majorant de \(A\).

\begin{resultatbox}
L'ensemble \(A\) est non vide et majoré.
\end{resultatbox}

% ------------------------------------------------------------

\subsubsection*{2. Montrer que \(1\leq r\leq 2\)}

On a défini
\[
r=\sup A.
\]

D'abord,
\[
1\in\R_+
\qquad\text{et}\qquad
1^2=1\leq 2,
\]
donc
\[
1\in A.
\]
Comme \(r\) est un majorant de \(A\), tout élément de \(A\) est inférieur ou égal à \(r\). En particulier,
\[
1\leq r.
\]

D'autre part, on a montré à la question précédente que \(2\) est un majorant de \(A\). Or \(r=\sup A\) est le plus petit des majorants de \(A\). Donc
\[
r\leq 2.
\]

\begin{resultatbox}
On a bien
\[
1\leq r\leq 2.
\]
\end{resultatbox}

% ------------------------------------------------------------

\subsubsection*{3. Montrer que l'on ne peut pas avoir \(r^2<2\)}

On raisonne par l'absurde.

Supposons que
\[
r^2<2.
\]
Alors
\[
2-r^2>0.
\]
On souhaite construire un réel strictement plus grand que \(r\), encore dans \(A\). Pour cela, on choisit un réel \(\eta>0\) assez petit pour que
\[
(r+\eta)^2<2.
\]

On développe :
\[
(r+\eta)^2=r^2+2r\eta+\eta^2.
\]
Comme \(1\leq r\leq 2\), on a en particulier \(r\leq 2\). Prenons par exemple
\[
0<\eta\leq 1.
\]
Alors
\[
2r\eta+\eta^2\leq 4\eta+\eta=5\eta.
\]
Il suffit donc d'imposer
\[
5\eta<2-r^2.
\]
On peut choisir
\[
0<\eta<\min\left(1,\frac{2-r^2}{5}\right).
\]
Alors
\[
(r+\eta)^2
=
r^2+2r\eta+\eta^2
\leq r^2+5\eta
<r^2+(2-r^2)=2.
\]
Donc
\[
r+\eta\in A.
\]
Mais
\[
r+\eta>r.
\]
Cela contredit le fait que \(r\) est un majorant de \(A\).

Donc l'hypothèse \(r^2<2\) est impossible.

\begin{resultatbox}
On ne peut pas avoir \(r^2<2\).
\end{resultatbox}

% ------------------------------------------------------------

\subsubsection*{4. Montrer que l'on ne peut pas avoir \(r^2>2\)}

On raisonne encore par l'absurde.

Supposons que
\[
r^2>2.
\]
Alors
\[
r^2-2>0.
\]
On veut montrer qu'un réel strictement plus petit que \(r\) serait encore un majorant de \(A\), ce qui contredirait la minimalité de \(r=\sup A\).

Comme \(r\geq 1\), on peut choisir \(\eta>0\) tel que
\[
0<\eta<r
\]
et
\[
(r-\eta)^2>2.
\]
En effet,
\[
(r-\eta)^2=r^2-2r\eta+\eta^2.
\]
Comme \(\eta^2\geq 0\), on a
\[
(r-\eta)^2\geq r^2-2r\eta.
\]
Il suffit donc de choisir \(\eta\) tel que
\[
r^2-2r\eta>2.
\]
Cela équivaut à
\[
2r\eta<r^2-2.
\]
On peut donc prendre
\[
0<\eta<\min\left(r,\frac{r^2-2}{2r}\right).
\]
Alors
\[
(r-\eta)^2>2.
\]

Montrons que \(r-\eta\) est un majorant de \(A\). Soit \(x\in A\). Alors
\[
x\geq 0
\qquad\text{et}\qquad
x^2\leq 2.
\]
Si l'on avait
\[
x>r-\eta,
\]
comme \(x\geq 0\) et \(r-\eta>0\), on aurait
\[
x^2>(r-\eta)^2>2,
\]
ce qui contredit \(x^2\leq 2\). Donc
\[
x\leq r-\eta.
\]
Ainsi
\[
\forall x\in A,\qquad x\leq r-\eta.
\]
Donc \(r-\eta\) est un majorant de \(A\).

Mais
\[
r-\eta<r.
\]
Cela contredit le fait que \(r\) est le plus petit majorant de \(A\).

\begin{resultatbox}
On ne peut pas avoir \(r^2>2\).
\end{resultatbox}

% ------------------------------------------------------------

\subsubsection*{5. Conclure que \(r^2=2\)}

D'après les deux questions précédentes,
\[
r^2<2
\]
est impossible, et
\[
r^2>2
\]
est impossible.

Comme \(r^2\) et \(2\) sont deux réels comparables, il reste nécessairement
\[
r^2=2.
\]

\begin{resultatbox}
On a
\[
r^2=2.
\]
\end{resultatbox}

% ------------------------------------------------------------

\subsubsection*{6. Montrer que \(r\) est l'unique réel positif vérifiant \(x^2=2\)}

On sait déjà que
\[
r\geq 0
\qquad\text{et}\qquad
r^2=2.
\]

Soit \(x\in\R_+\) tel que
\[
x^2=2.
\]
Alors
\[
x^2-r^2=0.
\]
Donc
\[
(x-r)(x+r)=0.
\]
Or \(x\geq 0\) et \(r\geq 0\), donc
\[
x+r\geq 0.
\]
De plus, \(r^2=2\), donc \(r>0\), et \(x^2=2\), donc \(x>0\). Ainsi
\[
x+r>0.
\]
Par conséquent,
\[
x-r=0.
\]
Donc
\[
x=r.
\]

\begin{resultatbox}
Le réel \(r\) est l'unique réel positif dont le carré vaut \(2\). On le note
\[
r=\sqrt2.
\]
\end{resultatbox}

% ============================================================

\subsection*{Partie B -- Une fonction itérative}

On définit
\[
\varphi(x)=\frac12\left(x+\frac2x\right)
\]
sur \(]0,+\infty[\).

% ------------------------------------------------------------

\subsubsection*{7. Montrer que \(\varphi\) est bien définie et continue sur \(]0,+\infty[\)}

Pour \(x>0\), on a \(x\neq 0\). Donc l'expression
\[
\frac2x
\]
est bien définie.

De plus,
\[
x\mapsto x
\]
est continue sur \(]0,+\infty[\), et
\[
x\mapsto \frac2x
\]
est continue sur \(]0,+\infty[\). Par somme et multiplication par une constante, la fonction
\[
\varphi:x\mapsto \frac12\left(x+\frac2x\right)
\]
est continue sur \(]0,+\infty[\).

Enfin, si \(x>0\), alors
\[
x>0
\qquad\text{et}\qquad
\frac2x>0.
\]
Donc
\[
\varphi(x)>0.
\]
Ainsi \(\varphi\) envoie bien \(]0,+\infty[\) dans \(]0,+\infty[\).

\begin{resultatbox}
La fonction \(\varphi\) est bien définie et continue sur \(]0,+\infty[\).
\end{resultatbox}

% ------------------------------------------------------------

\subsubsection*{8. Montrer l'identité fondamentale}

On veut montrer que, pour tout \(x>0\),
\[
\varphi(x)-\sqrt2
=
\frac{(x-\sqrt2)^2}{2x}.
\]

Calculons :
\[
\varphi(x)-\sqrt2
=
\frac12\left(x+\frac2x\right)-\sqrt2.
\]
On met au même dénominateur :
\[
\varphi(x)-\sqrt2
=
\frac{x}{2}+\frac1x-\sqrt2
=
\frac{x^2+2-2\sqrt2 x}{2x}.
\]
Or
\[
x^2-2\sqrt2 x+2=(x-\sqrt2)^2.
\]
Donc
\[
\varphi(x)-\sqrt2
=
\frac{(x-\sqrt2)^2}{2x}.
\]

\begin{resultatbox}
Pour tout \(x>0\),
\[
\varphi(x)-\sqrt2
=
\frac{(x-\sqrt2)^2}{2x}.
\]
\end{resultatbox}

% ------------------------------------------------------------

\subsubsection*{9. En déduire que \(\varphi(x)\geq\sqrt2\)}

Pour \(x>0\), on a
\[
2x>0
\]
et
\[
(x-\sqrt2)^2\geq 0.
\]
Donc
\[
\frac{(x-\sqrt2)^2}{2x}\geq 0.
\]
D'après l'identité précédente,
\[
\varphi(x)-\sqrt2\geq 0.
\]
Ainsi
\[
\varphi(x)\geq \sqrt2.
\]

\begin{resultatbox}
Pour tout \(x>0\),
\[
\varphi(x)\geq \sqrt2.
\]
\end{resultatbox}

% ------------------------------------------------------------

\subsubsection*{10. Montrer que \(\varphi(x)\leq x\) pour \(x\geq\sqrt2\)}

Soit \(x\geq\sqrt2\). Comme \(x>0\),
\[
\varphi(x)-x
=
\frac12\left(x+\frac2x\right)-x
=
\frac{2-x^2}{2x}.
\]
Or \(x\geq\sqrt2\), donc
\[
x^2\geq 2.
\]
Ainsi
\[
2-x^2\leq 0.
\]
Comme
\[
2x>0,
\]
on obtient
\[
\frac{2-x^2}{2x}\leq 0.
\]
Donc
\[
\varphi(x)-x\leq 0.
\]
Autrement dit,
\[
\varphi(x)\leq x.
\]

\begin{resultatbox}
Pour tout \(x\geq\sqrt2\),
\[
\varphi(x)\leq x.
\]
\end{resultatbox}

% ------------------------------------------------------------

\subsubsection*{11. Stabilité de \([\sqrt2,+\infty[\)}

Soit
\[
x\in[\sqrt2,+\infty[.
\]
Alors \(x>0\). D'après la question 9,
\[
\varphi(x)\geq\sqrt2.
\]
Donc
\[
\varphi(x)\in[\sqrt2,+\infty[.
\]

\begin{resultatbox}
L'intervalle \([\sqrt2,+\infty[\) est stable par \(\varphi\).
\end{resultatbox}

% ============================================================

\subsection*{Partie C -- Convergence de la suite de Héron}

On définit
\[
u_0=2,
\qquad
u_{n+1}=\varphi(u_n).
\]

% ------------------------------------------------------------

\subsubsection*{12. Montrer que \(u_n\geq\sqrt2\)}

On procède par récurrence.

Pour \(n=0\),
\[
u_0=2.
\]
Or
\[
2\geq\sqrt2,
\]
donc la propriété est vraie au rang \(0\).

Supposons maintenant que, pour un certain \(n\in\N\),
\[
u_n\geq\sqrt2.
\]
Alors
\[
u_n\in[\sqrt2,+\infty[.
\]
Comme cet intervalle est stable par \(\varphi\), on obtient
\[
u_{n+1}=\varphi(u_n)\in[\sqrt2,+\infty[.
\]
Donc
\[
u_{n+1}\geq\sqrt2.
\]

Par récurrence,
\[
\forall n\in\N,\qquad u_n\geq\sqrt2.
\]

% ------------------------------------------------------------

\subsubsection*{13. Montrer que \((u_n)\) est décroissante}

D'après la question précédente,
\[
u_n\geq\sqrt2
\]
pour tout \(n\in\N\). D'après la question 10, pour tout \(x\geq\sqrt2\),
\[
\varphi(x)\leq x.
\]
En appliquant cela à \(x=u_n\), on obtient
\[
u_{n+1}=\varphi(u_n)\leq u_n.
\]
Donc
\[
\forall n\in\N,\qquad u_{n+1}\leq u_n.
\]

\begin{resultatbox}
La suite \((u_n)\) est décroissante.
\end{resultatbox}

% ------------------------------------------------------------

\subsubsection*{14. En déduire que \((u_n)\) converge}

La suite \((u_n)\) est décroissante et minorée par \(\sqrt2\). En effet,
\[
\forall n\in\N,\qquad u_n\geq\sqrt2.
\]
D'après le théorème de convergence des suites monotones, toute suite décroissante et minorée converge.

Donc il existe un réel \(\ell\) tel que
\[
u_n\longrightarrow \ell.
\]
De plus, comme \(u_n\geq\sqrt2\) pour tout \(n\), le passage à la limite dans les inégalités donne
\[
\ell\geq\sqrt2>0.
\]

% ------------------------------------------------------------

\subsubsection*{15. Montrer que \(\ell=\sqrt2\)}

On sait que
\[
u_{n+1}=\varphi(u_n).
\]
Comme
\[
u_n\longrightarrow \ell
\]
et que \(\varphi\) est continue sur \(]0,+\infty[\), avec \(\ell>0\), on a
\[
\varphi(u_n)\longrightarrow \varphi(\ell).
\]
Or
\[
u_{n+1}\longrightarrow \ell
\]
puisque \((u_n)\) et sa suite décalée ont la même limite.

Ainsi, en passant à la limite dans
\[
u_{n+1}=\varphi(u_n),
\]
on obtient
\[
\ell=\varphi(\ell).
\]
Donc
\[
\ell=\frac12\left(\ell+\frac2\ell\right).
\]
Comme \(\ell>0\), on peut multiplier par \(2\ell\) :
\[
2\ell^2=\ell^2+2.
\]
Donc
\[
\ell^2=2.
\]
Comme \(\ell>0\), par unicité du réel positif de carré \(2\),
\[
\ell=\sqrt2.
\]

\begin{resultatbox}
La suite \((u_n)\) converge vers \(\sqrt2\).
\end{resultatbox}

% ------------------------------------------------------------

\subsubsection*{16. Inégalité d'erreur}

D'après l'identité de la question 8,
\[
u_{n+1}-\sqrt2
=
\frac{(u_n-\sqrt2)^2}{2u_n}.
\]
Comme \(u_n\geq\sqrt2>0\), on a
\[
2u_n\geq 2\sqrt2.
\]
En particulier,
\[
2u_n\geq 2.
\]
Donc
\[
\frac1{2u_n}\leq \frac12.
\]
Ainsi
\[
u_{n+1}-\sqrt2
=
\frac{(u_n-\sqrt2)^2}{2u_n}
\leq
\frac12 (u_n-\sqrt2)^2.
\]
De plus,
\[
u_{n+1}\geq \sqrt2,
\]
donc
\[
u_{n+1}-\sqrt2\geq 0.
\]

\begin{resultatbox}
Pour tout \(n\in\N\),
\[
0\leq u_{n+1}-\sqrt2
\leq
\frac12 (u_n-\sqrt2)^2.
\]
\end{resultatbox}

% ------------------------------------------------------------

\subsubsection*{17. Interprétation qualitative}

L'inégalité
\[
0\leq u_{n+1}-\sqrt2
\leq
\frac12 (u_n-\sqrt2)^2
\]
montre que l'erreur au rang \(n+1\) est contrôlée par le carré de l'erreur au rang \(n\).

Cela signifie que lorsque l'erreur devient petite, son carré devient beaucoup plus petit. Par exemple, si
\[
u_n-\sqrt2\approx 10^{-2},
\]
alors
\[
u_{n+1}-\sqrt2
\]
est majorée par une quantité de l'ordre de
\[
10^{-4}.
\]
Cela rend plausible la rapidité de convergence de la méthode.

\begin{remarquebox}
On ne parle pas ici de convergence quadratique au sens avancé du terme, car cette notion n'est pas nécessaire au programme. On se contente d'une interprétation qualitative à partir de l'inégalité obtenue.
\end{remarquebox}

% ============================================================

\subsection*{Partie D -- Deux suites adjacentes cachées}

On pose
\[
v_n=\frac2{u_n}.
\]

% ------------------------------------------------------------

\subsubsection*{18. Montrer que \(0<v_n\leq\sqrt2\leq u_n\)}

On sait que
\[
u_n\geq\sqrt2>0.
\]
Donc
\[
v_n=\frac2{u_n}>0.
\]
De plus, comme \(u_n\geq\sqrt2>0\), en prenant les inverses, on obtient
\[
\frac1{u_n}\leq \frac1{\sqrt2}.
\]
Donc
\[
\frac2{u_n}\leq \frac2{\sqrt2}.
\]
Or
\[
\frac2{\sqrt2}=\sqrt2.
\]
Ainsi
\[
v_n\leq\sqrt2.
\]
On a déjà
\[
u_n\geq\sqrt2.
\]
Donc
\[
0<v_n\leq\sqrt2\leq u_n.
\]

% ------------------------------------------------------------

\subsubsection*{19. Montrer que \((v_n)\) est croissante}

On sait que \((u_n)\) est décroissante et que
\[
u_n>0
\]
pour tout \(n\). Donc
\[
u_{n+1}\leq u_n.
\]
En prenant les inverses, puisque les deux termes sont strictement positifs,
\[
\frac1{u_{n+1}}\geq \frac1{u_n}.
\]
En multipliant par \(2>0\), on obtient
\[
\frac2{u_{n+1}}\geq \frac2{u_n}.
\]
Donc
\[
v_{n+1}\geq v_n.
\]
Ainsi, \((v_n)\) est croissante.

% ------------------------------------------------------------

\subsubsection*{20. Montrer que \(u_n-v_n\to 0\)}

On a
\[
u_n\longrightarrow \sqrt2.
\]
Comme la fonction
\[
x\mapsto \frac2x
\]
est continue en \(\sqrt2>0\), on obtient
\[
v_n=\frac2{u_n}\longrightarrow \frac2{\sqrt2}=\sqrt2.
\]
Donc
\[
u_n-v_n\longrightarrow \sqrt2-\sqrt2=0.
\]

% ------------------------------------------------------------

\subsubsection*{21. \((u_n)\) et \((v_n)\) même limite?}

On rappelle que deux suites \((a_n)\) et \((b_n)\) sont adjacentes si l'une est croissante, l'autre décroissante, et si leur différence tend vers \(0\).

Ici :
\[
(v_n)\ \text{est croissante},
\]
\[
(u_n)\ \text{est décroissante},
\]
et
\[
u_n-v_n\longrightarrow 0.
\]
De plus,
\[
v_n\leq u_n.
\]
Donc les suites \((v_n)\) et \((u_n)\) sont adjacentes. elles ont donc une même limite

% ------------------------------------------------------------

\subsubsection*{22. Limite commune}

Deux suites adjacentes convergent vers la même limite. D'après la question précédente, \((u_n)\) et \((v_n)\) sont adjacentes, donc elles ont une limite commune.

On sait déjà que
\[
u_n\longrightarrow \sqrt2.
\]
Donc
\[
v_n\longrightarrow \sqrt2.
\]

\begin{resultatbox}
Les deux suites \((u_n)\) et \((v_n)\) convergent vers \(\sqrt2\).
\end{resultatbox}

\newpage

% ============================================================
% PROBLEME II
% ============================================================

\section*{Problème II -- Une fonction prolongée, une bijection et des inégalités par accroissements finis}

On considère la fonction \(f\) définie sur \(]-1,+\infty[\) par
\[
f(x)=
\begin{cases}
\dfrac{\sqrt{1+x}-1}{x} & \text{si } x\neq 0,\\[1.2em]
\dfrac12 & \text{si } x=0.
\end{cases}
\]

% ============================================================

\subsection*{Partie A -- Prolongement et premières propriétés}

\subsubsection*{1. Rationalisation}

Pour \(x\neq 0\), on a
\[
f(x)=\frac{\sqrt{1+x}-1}{x}.
\]
On multiplie le numérateur et le dénominateur par la quantité conjuguée
\[
\sqrt{1+x}+1.
\]
Alors
\[
f(x)
=
\frac{(\sqrt{1+x}-1)(\sqrt{1+x}+1)}{x(\sqrt{1+x}+1)}.
\]
Or
\[
(\sqrt{1+x}-1)(\sqrt{1+x}+1)
=
(1+x)-1=x.
\]
Donc
\[
f(x)=\frac{x}{x(\sqrt{1+x}+1)}.
\]
Comme \(x\neq 0\),
\[
f(x)=\frac1{\sqrt{1+x}+1}.
\]

\begin{resultatbox}
Pour tout \(x\in]-1,+\infty[\setminus\{0\}\),
\[
f(x)=\frac1{\sqrt{1+x}+1}.
\]
\end{resultatbox}

% ------------------------------------------------------------

\subsubsection*{2. Continuité en \(0\)}

Pour \(x\neq 0\),
\[
f(x)=\frac1{\sqrt{1+x}+1}.
\]
Lorsque \(x\to 0\),
\[
\sqrt{1+x}\to 1.
\]
Donc
\[
\sqrt{1+x}+1\to 2,
\]
et par quotient,
\[
f(x)\to \frac12.
\]
Or
\[
f(0)=\frac12.
\]
Donc
\[
\lim_{x\to 0}f(x)=f(0).
\]

\begin{resultatbox}
La fonction \(f\) est continue en \(0\).
\end{resultatbox}

% ------------------------------------------------------------

\subsubsection*{3. Continuité sur \(]-1,+\infty[\)}

Sur \(]-1,+\infty[\setminus\{0\}\), la fonction
\[
x\mapsto \frac1{\sqrt{1+x}+1}
\]
est continue car elle est obtenue par composition, somme et quotient de fonctions continues, et le dénominateur ne s'annule jamais.

On vient de montrer que \(f\) est continue en \(0\). Donc \(f\) est continue en tout point de \(]-1,+\infty[\).

\begin{resultatbox}
La fonction \(f\) est continue sur \(]-1,+\infty[\).
\end{resultatbox}

% ------------------------------------------------------------

\subsubsection*{4. Limites aux bornes}

Pour tout \(x\in]-1,+\infty[\),
\[
f(x)=\frac1{\sqrt{1+x}+1},
\]
cette formule étant valable aussi en \(0\).

Lorsque
\[
x\to -1^+,
\]
on a
\[
1+x\to 0^+
\]
donc
\[
\sqrt{1+x}\to 0.
\]
Ainsi
\[
f(x)\to \frac1{0+1}=1.
\]
Donc
\[
\lim_{x\to -1^+}f(x)=1.
\]

Lorsque
\[
x\to+\infty,
\]
on a
\[
\sqrt{1+x}\to+\infty,
\]
donc
\[
\sqrt{1+x}+1\to+\infty.
\]
Ainsi
\[
f(x)=\frac1{\sqrt{1+x}+1}\to 0.
\]
Donc
\[
\lim_{x\to+\infty}f(x)=0.
\]

\begin{resultatbox}
On a
\[
\lim_{x\to -1^+}f(x)=1,
\qquad
\lim_{x\to+\infty}f(x)=0.
\]
\end{resultatbox}

% ------------------------------------------------------------

\subsubsection*{5. Montrer que \(0<f(x)<1\)}

Soit \(x\in]-1,+\infty[\). Alors
\[
1+x>0,
\]
donc
\[
\sqrt{1+x}>0.
\]
Ainsi
\[
\sqrt{1+x}+1>1.
\]
Par conséquent,
\[
0<\frac1{\sqrt{1+x}+1}<1.
\]
Donc
\[
0<f(x)<1.
\]

% ============================================================

\subsection*{Partie B -- Dérivabilité et variations}

\subsubsection*{6. Dérivabilité et calcul de \(f'(x)\)}

Grâce à la formule rationalisée, on peut écrire pour tout \(x\in]-1,+\infty[\),
\[
f(x)=\frac1{\sqrt{1+x}+1}.
\]
Cette expression est dérivable sur \(]-1,+\infty[\), car \(x\mapsto \sqrt{1+x}\) est dérivable sur cet intervalle.

Posons
\[
u(x)=\sqrt{1+x}+1.
\]
Alors
\[
f(x)=\frac1{u(x)}.
\]
On a
\[
u'(x)=\frac1{2\sqrt{1+x}}.
\]
Donc
\[
f'(x)=-\frac{u'(x)}{u(x)^2}.
\]
Ainsi
\[
f'(x)
=
-\frac{1}{2\sqrt{1+x}(\sqrt{1+x}+1)^2}.
\]

\begin{resultatbox}
Pour tout \(x>-1\),
\[
f'(x)
=
-\frac{1}{2\sqrt{1+x}(\sqrt{1+x}+1)^2}.
\]
\end{resultatbox}

\begin{remarquebox}
Cette formule est valable en particulier en \(0\). On a donc évité un calcul direct du taux d'accroissement en \(0\) en utilisant le prolongement rationalisé.
\end{remarquebox}

% ------------------------------------------------------------

\subsubsection*{7. Montrer que \(f\) est strictement décroissante}

Pour tout \(x>-1\),
\[
\sqrt{1+x}>0
\]
et
\[
(\sqrt{1+x}+1)^2>0.
\]
Donc
\[
2\sqrt{1+x}(\sqrt{1+x}+1)^2>0.
\]
Par conséquent,
\[
f'(x)<0.
\]
Ainsi \(f\) est strictement décroissante sur \(]-1,+\infty[\).

% ------------------------------------------------------------

\subsubsection*{8. Retrouver la stricte décroissance par Rolle}

On veut prouver la stricte décroissance sans se contenter de dire que \(f'<0\).

Soient \(x_1,x_2\in]-1,+\infty[\) tels que
\[
x_1<x_2.
\]
On veut montrer
\[
f(x_1)>f(x_2).
\]

Supposons par l'absurde que
\[
f(x_1)\leq f(x_2).
\]
Comme \(f'(x)<0\) pour tout \(x\), la fonction \(f\) ne peut pas avoir deux valeurs égales en deux points distincts. Justifions-le par Rolle.

Si l'on avait
\[
f(x_1)=f(x_2),
\]
alors \(f\) serait continue sur \([x_1,x_2]\), dérivable sur \(]x_1,x_2[\), et le théorème de Rolle donnerait un réel \(c\in]x_1,x_2[\) tel que
\[
f'(c)=0.
\]
C'est impossible puisque
\[
f'(c)<0.
\]
Donc \(f\) est injective.

Pour obtenir la décroissance stricte, on peut aussi utiliser le théorème des accroissements finis. Il existe \(c\in]x_1,x_2[\) tel que
\[
f(x_2)-f(x_1)=f'(c)(x_2-x_1).
\]
Comme
\[
f'(c)<0
\]
et
\[
x_2-x_1>0,
\]
on obtient
\[
f(x_2)-f(x_1)<0.
\]
Donc
\[
f(x_2)<f(x_1).
\]

\begin{resultatbox}
La fonction \(f\) est strictement décroissante sur \(]-1,+\infty[\).
\end{resultatbox}

% ------------------------------------------------------------

\subsubsection*{9. Bijection de \(]-1,+\infty[\) sur \(]0,1[\)}

On sait que \(f\) est continue et strictement décroissante sur l'intervalle \(]-1,+\infty[\).

De plus,
\[
\lim_{x\to -1^+}f(x)=1
\]
et
\[
\lim_{x\to+\infty}f(x)=0.
\]
Comme \(f\) est continue et strictement monotone, elle réalise une bijection de \(]-1,+\infty[\) sur son image.

Son image est l'intervalle ouvert formé par les valeurs strictement comprises entre les limites aux bornes :
\[
f(]-1,+\infty[)=]0,1[.
\]

\begin{resultatbox}
La fonction \(f\) réalise une bijection de \(]-1,+\infty[\) sur \(]0,1[\).
\end{resultatbox}

% ------------------------------------------------------------

\subsubsection*{10. Déterminer \(f^{-1}\)}

Soit
\[
y\in]0,1[.
\]
On cherche \(x>-1\) tel que
\[
y=f(x).
\]
D'après la forme rationalisée,
\[
y=\frac1{\sqrt{1+x}+1}.
\]
Comme \(y>0\),
\[
\sqrt{1+x}+1=\frac1y.
\]
Donc
\[
\sqrt{1+x}=\frac1y-1.
\]
Comme \(y\in]0,1[\), on a
\[
\frac1y>1,
\]
donc
\[
\frac1y-1>0.
\]
On peut donc élever au carré :
\[
1+x=\left(\frac1y-1\right)^2.
\]
Ainsi
\[
x=\left(\frac1y-1\right)^2-1.
\]

\begin{resultatbox}
Pour tout \(y\in]0,1[\),
\[
f^{-1}(y)=\left(\frac1y-1\right)^2-1.
\]
\end{resultatbox}

% ============================================================

\subsection*{Partie C -- Contrôle de la réciproque}

On note
\[
g=f^{-1}.
\]
Donc
\[
g(y)=\left(\frac1y-1\right)^2-1
\]
sur \(]0,1[\).

% ------------------------------------------------------------

\subsubsection*{11. Domaine et expression de \(g\)}

D'après la question précédente,
\[
g: ]0,1[\longrightarrow ]-1,+\infty[
\]
est définie par
\[
g(y)=\left(\frac1y-1\right)^2-1.
\]
On peut simplifier :
\[
\frac1y-1=\frac{1-y}{y}.
\]
Donc
\[
g(y)=\frac{(1-y)^2}{y^2}-1.
\]
En développant,
\[
g(y)=\frac{1-2y+y^2-y^2}{y^2}
=
\frac{1-2y}{y^2}.
\]
Ainsi
\[
g(y)=\frac1{y^2}-\frac2y.
\]

% ------------------------------------------------------------

\subsubsection*{12. Dérivabilité et calcul de \(g'(y)\)}

La fonction \(g\) est dérivable sur \(]0,1[\), car elle est une combinaison de puissances de \(y\) avec \(y\neq0\).

Avec
\[
g(y)=y^{-2}-2y^{-1},
\]
on obtient
\[
g'(y)=-2y^{-3}+2y^{-2}.
\]
Donc
\[
g'(y)=\frac{-2}{y^3}+\frac2{y^2}
=
\frac{2y-2}{y^3}
=
\frac{2(y-1)}{y^3}.
\]

\begin{resultatbox}
Pour tout \(y\in]0,1[\),
\[
g'(y)=\frac{2(y-1)}{y^3}.
\]
\end{resultatbox}

% ------------------------------------------------------------

\subsubsection*{13. Contrôle sur \([a,b]\)}

Soient
\[
0<a<b<1.
\]
La fonction \(g\) est dérivable sur \(]0,1[\), donc elle est continue sur \([a,b]\) et dérivable sur \(]a,b[\).

La fonction \(g'\) est continue sur \([a,b]\), donc elle est bornée sur \([a,b]\). Il existe donc
\[
M_{a,b}>0
\]
tel que
\[
\forall t\in[a,b],
\qquad
|g'(t)|\leq M_{a,b}.
\]

Soient \(y,z\in[a,b]\). Si \(y=z\), l'inégalité est évidente.

Si \(y\neq z\), le théorème des accroissements finis appliqué à \(g\) entre \(y\) et \(z\) donne un réel \(c\) compris entre \(y\) et \(z\), donc \(c\in[a,b]\), tel que
\[
g(y)-g(z)=g'(c)(y-z).
\]
En prenant les valeurs absolues,
\[
|g(y)-g(z)|
=
|g'(c)|\,|y-z|
\leq
M_{a,b}|y-z|.
\]

\begin{resultatbox}
Il existe bien \(M_{a,b}>0\) tel que
\[
\forall y,z\in[a,b],
\qquad
|g(y)-g(z)|\leq M_{a,b}|y-z|.
\]
\end{resultatbox}

% ------------------------------------------------------------

\subsubsection*{14. Donner une valeur possible de \(M_{a,b}\)}

On a
\[
g'(y)=\frac{2(y-1)}{y^3}.
\]
Sur \(]0,1[\),
\[
y-1<0,
\]
donc
\[
|g'(y)|=\frac{2(1-y)}{y^3}.
\]
Si \(y\in[a,b]\), alors
\[
1-y\leq 1-a
\]
et
\[
y^3\geq a^3.
\]
Donc
\[
|g'(y)|
=
\frac{2(1-y)}{y^3}
\leq
\frac{2(1-a)}{a^3}.
\]
On peut donc prendre
\[
M_{a,b}=\frac{2(1-a)}{a^3}.
\]

\begin{remarquebox}
Cette constante n'est pas forcément optimale, mais elle est suffisante. Le but est de montrer l'existence d'une constante de contrôle.
\end{remarquebox}

% ------------------------------------------------------------

\subsubsection*{15. Impossibilité d'une constante globale sur \(]0,1[\)}

Supposons qu'il existe \(M>0\) tel que
\[
\forall y,z\in]0,1[,
\qquad
|g(y)-g(z)|\leq M|y-z|.
\]
Alors, fixons \(z=\frac12\) par exemple, ou raisonnons directement avec la dérivée.

Une telle inégalité impliquerait que tous les taux d'accroissement de \(g\) sont bornés par \(M\). En particulier, pour tout \(y\in]0,1[\), si \(g\) est dérivable en \(y\), on aurait
\[
|g'(y)|\leq M.
\]
Or
\[
|g'(y)|=\frac{2(1-y)}{y^3}.
\]
Quand
\[
y\to0^+,
\]
on a
\[
1-y\to1
\]
et
\[
y^3\to0^+.
\]
Donc
\[
|g'(y)|\to+\infty.
\]
Il est donc impossible que
\[
|g'(y)|\leq M
\]
pour tout \(y\in]0,1[\).

\begin{resultatbox}
Il n'existe pas de constante globale \(M\) valable sur tout \(]0,1[\).
\end{resultatbox}

% ============================================================

\subsection*{Partie D -- Une inégalité fine sans développement limité}

\subsubsection*{16. Expression de \(f'\)}

Déjà obtenue à la question 6 :
\[
f'(x)
=
-\frac{1}{2\sqrt{1+x}(\sqrt{1+x}+1)^2}.
\]

% ------------------------------------------------------------

\subsubsection*{17. Montrer que \(0\leq \frac12-f(x)\leq \frac{x}{8}\) pour \(x\geq0\)}

Soit \(x\geq0\).

Comme \(f\) est décroissante sur \(]-1,+\infty[\), on a pour \(x\geq0\),
\[
f(x)\leq f(0)=\frac12.
\]
Donc
\[
0\leq \frac12-f(x).
\]

Montrons maintenant la majoration. Sur \([0,x]\), on a pour tout \(t\in[0,x]\),
\[
t\geq0.
\]
Donc
\[
\sqrt{1+t}\geq1.
\]
Ainsi
\[
2\sqrt{1+t}(\sqrt{1+t}+1)^2
\geq
2\cdot1\cdot(1+1)^2
=
8.
\]
Donc
\[
|f'(t)|
=
\frac1{2\sqrt{1+t}(\sqrt{1+t}+1)^2}
\leq
\frac18.
\]
Par l'inégalité des accroissements finis sur \([0,x]\),
\[
|f(x)-f(0)|\leq \frac18 |x-0|.
\]
Donc
\[
\left|f(x)-\frac12\right|\leq \frac{x}{8}.
\]
Comme \(f(x)\leq1/2\), cela devient
\[
0\leq \frac12-f(x)\leq \frac{x}{8}.
\]

\begin{resultatbox}
Pour tout \(x\geq0\),
\[
0\leq \frac12-f(x)\leq \frac{x}{8}.
\]
\end{resultatbox}

% ------------------------------------------------------------

\subsubsection*{18. Reformulation}

Pour \(x>0\),
\[
f(x)=\frac{\sqrt{1+x}-1}{x}.
\]
Donc l'inégalité précédente devient
\[
0\leq \frac12-\frac{\sqrt{1+x}-1}{x}\leq \frac{x}{8}.
\]

% ------------------------------------------------------------

\subsubsection*{19. Approximation contrôlée de \(\sqrt{1+x}\)}

Pour \(x>0\), on part de
\[
0\leq \frac12-\frac{\sqrt{1+x}-1}{x}\leq \frac{x}{8}.
\]
Comme \(x>0\), on peut multiplier toute l'inégalité par \(x\) :
\[
0\leq \frac{x}{2}-(\sqrt{1+x}-1)\leq \frac{x^2}{8}.
\]
Or
\[
\frac{x}{2}-(\sqrt{1+x}-1)
=
1+\frac{x}{2}-\sqrt{1+x}.
\]
Donc
\[
0\leq 1+\frac{x}{2}-\sqrt{1+x}\leq \frac{x^2}{8}.
\]

\begin{resultatbox}
Pour tout \(x>0\),
\[
0\leq 1+\frac{x}{2}-\sqrt{1+x}\leq \frac{x^2}{8}.
\]
\end{resultatbox}

Cela signifie que
\[
\sqrt{1+x}
\]
est approchée par
\[
1+\frac{x}{2}
\]
avec une erreur comprise entre \(0\) et \(x^2/8\).

% ------------------------------------------------------------

\subsubsection*{20. Dérivabilité en \(0\) par définition}

On considère
\[
q(x)=\sqrt{1+x}.
\]
On sait que
\[
q(0)=1.
\]
Pour \(x\neq0\),
\[
\frac{q(x)-q(0)}{x}
=
\frac{\sqrt{1+x}-1}{x}.
\]
En rationalisant,
\[
\frac{\sqrt{1+x}-1}{x}
=
\frac1{\sqrt{1+x}+1}.
\]
Quand
\[
x\to0,
\]
on obtient
\[
\frac1{\sqrt{1+x}+1}\to\frac12.
\]
Donc \(q\) est dérivable en \(0\), et
\[
q'(0)=\frac12.
\]
Par la caractérisation locale de la dérivabilité,
\[
q(0+x)=q(0)+xq'(0)+x\eps(x),
\]
avec
\[
\eps(x)\to0
\quad\text{quand }x\to0.
\]
Donc
\[
\sqrt{1+x}=1+\frac{x}{2}+x\eps(x),
\]
avec
\[
\eps(x)\to0.
\]

\begin{remarquebox}
La question 19 donne une information plus précise pour \(x>0\) :
\[
0\leq 1+\frac{x}{2}-\sqrt{1+x}\leq\frac{x^2}{8}.
\]
Elle contrôle donc le signe et la taille de l'erreur.  
La dérivabilité donne seulement une erreur de la forme \(x\eps(x)\), avec \(\eps(x)\to0\).
\end{remarquebox}

\newpage

% ============================================================
% PROBLEME III
% ============================================================

\section*{Problème III -- Dichotomie, théorème des valeurs intermédiaires et approximation certifiée d'une racine}

On étudie la procédure de dichotomie.

Soit \(h\) une fonction continue sur \([a,b]\), telle que
\[
h(a)\leq0\leq h(b).
\]
On définit deux suites \((a_n)\) et \((b_n)\) par
\[
a_0=a,\qquad b_0=b.
\]
Pour tout \(n\in\N\), on pose
\[
m_n=\frac{a_n+b_n}{2}.
\]
Si
\[
h(m_n)\leq0,
\]
on pose
\[
a_{n+1}=m_n,\qquad b_{n+1}=b_n.
\]
Sinon,
\[
a_{n+1}=a_n,\qquad b_{n+1}=m_n.
\]

% ============================================================

\subsection*{Partie A -- Une preuve constructive du TVI}

\subsubsection*{1. Montrer que \(a_n\leq b_n\)}

On procède par récurrence.

Au rang \(0\),
\[
a_0=a,\qquad b_0=b.
\]
Comme \(h\) est définie sur \([a,b]\), on suppose bien
\[
a\leq b.
\]
Donc
\[
a_0\leq b_0.
\]

Supposons maintenant
\[
a_n\leq b_n.
\]
Alors
\[
m_n=\frac{a_n+b_n}{2}
\]
vérifie
\[
a_n\leq m_n\leq b_n.
\]
Deux cas sont possibles.

Premier cas :
\[
h(m_n)\leq0.
\]
Alors
\[
a_{n+1}=m_n,\qquad b_{n+1}=b_n.
\]
Comme \(m_n\leq b_n\), on obtient
\[
a_{n+1}\leq b_{n+1}.
\]

Deuxième cas :
\[
h(m_n)>0.
\]
Alors
\[
a_{n+1}=a_n,\qquad b_{n+1}=m_n.
\]
Comme \(a_n\leq m_n\), on obtient encore
\[
a_{n+1}\leq b_{n+1}.
\]

Par récurrence,
\[
\forall n\in\N,\qquad a_n\leq b_n.
\]

% ------------------------------------------------------------

\subsubsection*{2. Conservation de l'encadrement de signe}

On montre par récurrence que
\[
h(a_n)\leq0\leq h(b_n).
\]

Au rang \(0\),
\[
a_0=a,\qquad b_0=b.
\]
L'hypothèse donne
\[
h(a)\leq0\leq h(b),
\]
donc
\[
h(a_0)\leq0\leq h(b_0).
\]

Supposons maintenant que
\[
h(a_n)\leq0\leq h(b_n).
\]
On pose
\[
m_n=\frac{a_n+b_n}{2}.
\]

Si
\[
h(m_n)\leq0,
\]
alors
\[
a_{n+1}=m_n,\qquad b_{n+1}=b_n.
\]
Donc
\[
h(a_{n+1})=h(m_n)\leq0
\]
et
\[
h(b_{n+1})=h(b_n)\geq0.
\]
Ainsi
\[
h(a_{n+1})\leq0\leq h(b_{n+1}).
\]

Si
\[
h(m_n)>0,
\]
alors
\[
a_{n+1}=a_n,\qquad b_{n+1}=m_n.
\]
Donc
\[
h(a_{n+1})=h(a_n)\leq0
\]
et
\[
h(b_{n+1})=h(m_n)>0.
\]
Donc
\[
h(a_{n+1})\leq0\leq h(b_{n+1}).
\]

Par récurrence,
\[
\forall n\in\N,\qquad h(a_n)\leq0\leq h(b_n).
\]

% ------------------------------------------------------------

\subsubsection*{3. Monotonie de \((a_n)\) et \((b_n)\)}

À chaque étape, on a
\[
a_n\leq m_n\leq b_n.
\]

Si
\[
h(m_n)\leq0,
\]
alors
\[
a_{n+1}=m_n\geq a_n
\]
et
\[
b_{n+1}=b_n.
\]

Si
\[
h(m_n)>0,
\]
alors
\[
a_{n+1}=a_n
\]
et
\[
b_{n+1}=m_n\leq b_n.
\]

Dans tous les cas,
\[
a_{n+1}\geq a_n
\]
et
\[
b_{n+1}\leq b_n.
\]

\begin{resultatbox}
La suite \((a_n)\) est croissante et la suite \((b_n)\) est décroissante.
\end{resultatbox}

% ------------------------------------------------------------

\subsubsection*{4. Relation sur la longueur des intervalles}

On montre que
\[
b_{n+1}-a_{n+1}=\frac12(b_n-a_n).
\]

Si
\[
h(m_n)\leq0,
\]
alors
\[
a_{n+1}=m_n,\qquad b_{n+1}=b_n.
\]
Donc
\[
b_{n+1}-a_{n+1}
=
b_n-m_n.
\]
Or
\[
m_n=\frac{a_n+b_n}{2}.
\]
Ainsi
\[
b_n-m_n
=
b_n-\frac{a_n+b_n}{2}
=
\frac{b_n-a_n}{2}.
\]

Si
\[
h(m_n)>0,
\]
alors
\[
a_{n+1}=a_n,\qquad b_{n+1}=m_n.
\]
Donc
\[
b_{n+1}-a_{n+1}
=
m_n-a_n
=
\frac{a_n+b_n}{2}-a_n
=
\frac{b_n-a_n}{2}.
\]

Dans tous les cas,
\[
b_{n+1}-a_{n+1}=\frac12(b_n-a_n).
\]

% ------------------------------------------------------------

\subsubsection*{5. Formule explicite de la longueur}

On a
\[
b_{n+1}-a_{n+1}=\frac12(b_n-a_n).
\]
La suite
\[
\ell_n=b_n-a_n
\]
est donc géométrique de raison \(1/2\), avec
\[
\ell_0=b_0-a_0=b-a.
\]
Ainsi
\[
\ell_n=\frac{b-a}{2^n}.
\]
Donc
\[
b_n-a_n=\frac{b-a}{2^n}.
\]

% ------------------------------------------------------------

\subsubsection*{6. Montrer que \((a_n)\) et \((b_n)\) sont adjacentes}

On a montré que :
\[
(a_n)\ \text{est croissante},
\]
\[
(b_n)\ \text{est décroissante},
\]
et
\[
b_n-a_n=\frac{b-a}{2^n}.
\]
Comme
\[
2^n\to+\infty,
\]
on a
\[
\frac{b-a}{2^n}\to0.
\]
Donc
\[
b_n-a_n\to0.
\]
Les suites \((a_n)\) et \((b_n)\) sont donc adjacentes.

\begin{resultatbox}
Les suites \((a_n)\) et \((b_n)\) sont adjacentes.
\end{resultatbox}

% ------------------------------------------------------------

\subsubsection*{7. Localisation de la limite commune}

Puisque \((a_n)\) et \((b_n)\) sont adjacentes, elles convergent vers une même limite. Notons-la
\[
\alpha.
\]
Comme
\[
a_0=a\leq a_n\leq b_n\leq b_0=b
\]
pour tout \(n\), en passant à la limite, on obtient
\[
a\leq\alpha\leq b.
\]
Donc
\[
\alpha\in[a,b].
\]

% ------------------------------------------------------------

\subsubsection*{8. Montrer que \(h(\alpha)=0\)}

Pour tout \(n\),
\[
h(a_n)\leq0\leq h(b_n).
\]
Comme
\[
a_n\to\alpha
\qquad\text{et}\qquad
b_n\to\alpha,
\]
et comme \(h\) est continue sur \([a,b]\), on a
\[
h(a_n)\to h(\alpha)
\qquad\text{et}\qquad
h(b_n)\to h(\alpha).
\]

À partir de
\[
h(a_n)\leq0,
\]
le passage à la limite donne
\[
h(\alpha)\leq0.
\]
À partir de
\[
0\leq h(b_n),
\]
le passage à la limite donne
\[
0\leq h(\alpha).
\]
Donc
\[
h(\alpha)\leq0\leq h(\alpha).
\]
Ainsi
\[
h(\alpha)=0.
\]

\begin{resultatbox}
La limite commune \(\alpha\) des deux suites vérifie
\[
h(\alpha)=0.
\]
\end{resultatbox}

% ------------------------------------------------------------

\subsubsection*{9. Interprétation constructive du TVI}

Le théorème des valeurs intermédiaires affirme qu'une fonction continue qui prend des valeurs de signes opposés sur un segment possède au moins un zéro.

La procédure de dichotomie ne se contente pas d'affirmer l'existence du zéro : elle construit une suite d'intervalles emboîtés
\[
[a_n,b_n]
\]
dont les longueurs tendent vers \(0\), et dans chacun desquels un changement de signe est conservé.

La limite commune des bornes est alors un zéro de \(h\). La preuve donne donc non seulement l'existence, mais aussi une méthode d'approximation du zéro.

% ============================================================

\subsection*{Partie B -- Cas d'une racine unique}

On suppose maintenant que l'équation
\[
h(x)=0
\]
admet une unique solution dans \([a,b]\). On la note \(\alpha\).

% ------------------------------------------------------------

\subsubsection*{10. Montrer que \(a_n\leq\alpha\leq b_n\)}

D'après la partie A, pour chaque \(n\),
\[
h(a_n)\leq0\leq h(b_n).
\]
La fonction \(h\) est continue sur \([a_n,b_n]\), car \([a_n,b_n]\subset[a,b]\).

Par le théorème des valeurs intermédiaires, il existe un zéro de \(h\) dans \([a_n,b_n]\).

Mais par hypothèse, le zéro de \(h\) dans \([a,b]\) est unique : c'est \(\alpha\). Donc
\[
\alpha\in[a_n,b_n].
\]
Autrement dit,
\[
a_n\leq\alpha\leq b_n.
\]

% ------------------------------------------------------------

\subsubsection*{11. Estimation de l'erreur du milieu}

On définit
\[
c_n=\frac{a_n+b_n}{2}.
\]
Comme
\[
\alpha\in[a_n,b_n],
\]
le point \(c_n\) est le milieu de l'intervalle contenant \(\alpha\). La distance maximale entre le milieu d'un intervalle et un point de cet intervalle est la demi-longueur de l'intervalle.

Donc
\[
|c_n-\alpha|\leq \frac{b_n-a_n}{2}.
\]

% ------------------------------------------------------------

\subsubsection*{12. Estimation explicite}

On sait que
\[
b_n-a_n=\frac{b-a}{2^n}.
\]
Donc
\[
|c_n-\alpha|\leq \frac{b_n-a_n}{2}
=
\frac{b-a}{2^{n+1}}.
\]

\begin{resultatbox}
Pour tout \(n\in\N\),
\[
|c_n-\alpha|\leq \frac{b-a}{2^{n+1}}.
\]
\end{resultatbox}

% ------------------------------------------------------------

\subsubsection*{13. Montrer que \(c_n\to\alpha\)}

On a
\[
0\leq |c_n-\alpha|\leq \frac{b-a}{2^{n+1}}.
\]
Or
\[
\frac{b-a}{2^{n+1}}\to0.
\]
Par le théorème des gendarmes,
\[
|c_n-\alpha|\to0.
\]
Donc
\[
c_n\to\alpha.
\]

% ------------------------------------------------------------

\subsubsection*{14. Précision \(\varepsilon\)}

Soit
\[
\varepsilon>0.
\]
Comme
\[
\frac{b-a}{2^{n+1}}\to0,
\]
il existe un rang \(N\in\N\) tel que, pour tout \(n\geq N\),
\[
\frac{b-a}{2^{n+1}}\leq \varepsilon.
\]
D'après l'estimation précédente,
\[
|c_n-\alpha|\leq \frac{b-a}{2^{n+1}}.
\]
Donc, pour tout \(n\geq N\),
\[
|c_n-\alpha|\leq\varepsilon.
\]

% ------------------------------------------------------------

\subsubsection*{15. Approximation certifiée}

La dichotomie fournit une approximation certifiée car, à chaque rang \(n\), on connaît un intervalle explicite
\[
[a_n,b_n]
\]
qui contient la racine \(\alpha\). En choisissant le milieu
\[
c_n=\frac{a_n+b_n}{2},
\]
on sait rigoureusement que l'erreur est au plus
\[
\frac{b_n-a_n}{2}.
\]
On ne donne donc pas seulement une valeur approchée : on donne aussi une borne garantie de l'erreur.

% ============================================================

\subsection*{Partie C -- Application : la racine plastique}

On considère
\[
p(x)=x^3-x-1
\]
sur \([1,2]\).

% ------------------------------------------------------------

\subsubsection*{16. Continuité de \(p\)}

La fonction \(p\) est une fonction polynomiale. Elle est donc continue sur \(\R\), en particulier sur \([1,2]\).

% ------------------------------------------------------------

\subsubsection*{17. Existence d'une solution}

On calcule :
\[
p(1)=1^3-1-1=-1.
\]
Donc
\[
p(1)<0.
\]
Puis
\[
p(2)=2^3-2-1=8-3=5.
\]
Donc
\[
p(2)>0.
\]

La fonction \(p\) est continue sur \([1,2]\), et
\[
p(1)\leq0\leq p(2).
\]
Par le théorème des valeurs intermédiaires, il existe au moins un réel
\[
\rho\in[1,2]
\]
tel que
\[
p(\rho)=0.
\]

% ------------------------------------------------------------

\subsubsection*{18. Montrer que \(p\) est strictement croissante sur \([1,2]\)}

La fonction \(p\) est dérivable sur \([1,2]\), et
\[
p'(x)=3x^2-1.
\]
Pour \(x\in[1,2]\),
\[
x^2\geq1,
\]
donc
\[
3x^2-1\geq2>0.
\]
Ainsi
\[
p'(x)>0
\]
sur \([1,2]\).

Donc \(p\) est strictement croissante sur \([1,2]\).

% ------------------------------------------------------------

\subsubsection*{19. Unicité de la racine}

On sait qu'il existe au moins une racine de \(p\) dans \([1,2]\).

Comme \(p\) est strictement croissante sur \([1,2]\), elle est injective sur \([1,2]\). Elle ne peut donc pas prendre la valeur \(0\) en deux points distincts.

Donc l'équation
\[
p(x)=0
\]
admet une unique solution dans \([1,2]\). On la note
\[
\rho.
\]

% ------------------------------------------------------------

\subsubsection*{20. Trois premiers intervalles de dichotomie}

On applique la dichotomie à \(p\) sur \([1,2]\).

On a
\[
[a_0,b_0]=[1,2].
\]

Premier milieu :
\[
m_0=\frac{1+2}{2}=\frac32.
\]
Calculons
\[
p\left(\frac32\right)
=
\left(\frac32\right)^3-\frac32-1
=
\frac{27}{8}-\frac{12}{8}-\frac{8}{8}
=
\frac7{8}>0.
\]
Comme \(p(m_0)>0\), on garde la moitié gauche :
\[
[a_1,b_1]=\left[1,\frac32\right].
\]

Deuxième milieu :
\[
m_1=\frac{1+\frac32}{2}
=
\frac{\frac52}{2}
=
\frac54.
\]
Calculons
\[
p\left(\frac54\right)
=
\left(\frac54\right)^3-\frac54-1
=
\frac{125}{64}-\frac{80}{64}-\frac{64}{64}
=
-\frac{19}{64}<0.
\]
Donc on garde la moitié droite :
\[
[a_2,b_2]=\left[\frac54,\frac32\right].
\]

Troisième milieu :
\[
m_2=\frac{\frac54+\frac32}{2}
=
\frac{\frac{10}{8}+\frac{12}{8}}{2}
=
\frac{\frac{22}{8}}{2}
=
\frac{11}{8}.
\]
Calculons
\[
p\left(\frac{11}{8}\right)
=
\left(\frac{11}{8}\right)^3-\frac{11}{8}-1.
\]
On met au dénominateur \(512\) :
\[
\left(\frac{11}{8}\right)^3=\frac{1331}{512},
\]
\[
\frac{11}{8}=\frac{704}{512},
\qquad
1=\frac{512}{512}.
\]
Donc
\[
p\left(\frac{11}{8}\right)
=
\frac{1331-704-512}{512}
=
\frac{115}{512}>0.
\]
Donc on garde la moitié gauche :
\[
[a_3,b_3]=\left[\frac54,\frac{11}{8}\right].
\]

\begin{resultatbox}
Les premiers intervalles sont
\[
[a_0,b_0]=[1,2],
\]
\[
[a_1,b_1]=\left[1,\frac32\right],
\]
\[
[a_2,b_2]=\left[\frac54,\frac32\right],
\]
\[
[a_3,b_3]=\left[\frac54,\frac{11}{8}\right].
\]
\end{resultatbox}

% ------------------------------------------------------------

\subsubsection*{21. Vérifier que \(\frac54\leq\rho\leq\frac{11}{8}\)}

D'après la question précédente,
\[
[a_3,b_3]=\left[\frac54,\frac{11}{8}\right].
\]
La racine unique \(\rho\) appartient à chaque intervalle de localisation. Donc
\[
\frac54\leq\rho\leq\frac{11}{8}.
\]

% ------------------------------------------------------------

\subsubsection*{22. Étape supplémentaire}

On part de
\[
[a_3,b_3]=\left[\frac54,\frac{11}{8}\right].
\]
Le milieu est
\[
m_3=
\frac{\frac54+\frac{11}{8}}{2}
=
\frac{\frac{10}{8}+\frac{11}{8}}{2}
=
\frac{\frac{21}{8}}{2}
=
\frac{21}{16}.
\]
Calculons
\[
p\left(\frac{21}{16}\right)
=
\left(\frac{21}{16}\right)^3-\frac{21}{16}-1.
\]
On met au dénominateur \(4096\) :
\[
\left(\frac{21}{16}\right)^3=\frac{9261}{4096},
\]
\[
\frac{21}{16}=\frac{5376}{4096},
\qquad
1=\frac{4096}{4096}.
\]
Donc
\[
p\left(\frac{21}{16}\right)
=
\frac{9261-5376-4096}{4096}
=
-\frac{211}{4096}<0.
\]
Comme
\[
p\left(\frac{21}{16}\right)<0
\]
et
\[
p\left(\frac{11}{8}\right)>0,
\]
on obtient le nouvel intervalle
\[
\left[\frac{21}{16},\frac{11}{8}\right].
\]
Donc
\[
\frac{21}{16}\leq\rho\leq\frac{11}{8}.
\]

% ------------------------------------------------------------

\subsubsection*{23. Encadrement par un intervalle de longueur inférieure ou égale à \(1/16\)}

L'intervalle obtenu à la question précédente est
\[
\left[\frac{21}{16},\frac{11}{8}\right].
\]
Sa longueur vaut
\[
\frac{11}{8}-\frac{21}{16}
=
\frac{22}{16}-\frac{21}{16}
=
\frac1{16}.
\]
Donc il s'agit bien d'un encadrement de \(\rho\) par un intervalle de longueur inférieure ou égale à \(1/16\).

\begin{resultatbox}
On a
\[
\rho\in\left[\frac{21}{16},\frac{11}{8}\right],
\]
et cet intervalle est de longueur \(1/16\).
\end{resultatbox}

% ------------------------------------------------------------

\subsubsection*{24. Nombre d'étapes pour garantir une approximation à \(10^{-3}\)}

Ici,
\[
a=1,\qquad b=2,
\]
donc
\[
b-a=1.
\]
L'erreur commise en prenant le milieu \(c_n\) de l'intervalle \([a_n,b_n]\) est majorée par
\[
|c_n-\rho|\leq \frac{b-a}{2^{n+1}}
=
\frac1{2^{n+1}}.
\]
On veut
\[
\frac1{2^{n+1}}\leq 10^{-3}.
\]
Cela équivaut à
\[
2^{n+1}\geq 1000.
\]
On sait que
\[
2^{10}=1024\geq1000.
\]
Il suffit donc de prendre
\[
n+1=10,
\]
c'est-à-dire
\[
n=9.
\]

\begin{resultatbox}
Après \(9\) étapes, le milieu de l'intervalle donne une approximation de \(\rho\) à \(10^{-3}\) près.
\end{resultatbox}

\begin{remarquebox}
Selon la convention utilisée pour compter les étapes, on peut dire : après avoir construit l'intervalle \([a_9,b_9]\), son milieu est à distance au plus \(10^{-3}\) de \(\rho\).
\end{remarquebox}

% ------------------------------------------------------------

\subsubsection*{25. Montrer que \(c_n\to\rho\)}

La question 13, appliquée ici, donne directement
\[
|c_n-\rho|\leq\frac1{2^{n+1}}.
\]
Comme
\[
\frac1{2^{n+1}}\to0,
\]
le théorème des gendarmes donne
\[
|c_n-\rho|\to0.
\]
Donc
\[
c_n\to\rho.
\]

% ============================================================

\subsection*{Partie D -- Existence, unicité et algorithme}

\subsubsection*{26. L'unicité était-elle nécessaire dans la partie A ?}

Non. Dans la partie A, on n'a jamais supposé que le zéro de \(h\) était unique.

On a seulement utilisé :
\[
h(a)\leq0\leq h(b),
\]
la continuité de \(h\), et la construction par dichotomie. Cela suffit pour construire deux suites adjacentes dont la limite commune est un zéro de \(h\).

L'unicité n'est donc pas nécessaire pour démontrer l'existence d'au moins un zéro.

% ------------------------------------------------------------

\subsubsection*{27. Où l'unicité est-elle utilisée dans la partie B ?}

L'unicité est utilisée pour affirmer que la racine \(\alpha\) donnée au départ appartient à tous les intervalles \([a_n,b_n]\).

Sans unicité, le théorème des valeurs intermédiaires montrerait seulement que chaque intervalle \([a_n,b_n]\) contient au moins un zéro de \(h\). Mais il ne garantirait pas qu'il contient un zéro particulier choisi à l'avance.

Ainsi, l'unicité permet d'écrire
\[
a_n\leq\alpha\leq b_n
\]
pour la même racine \(\alpha\) à tous les rangs.

% ------------------------------------------------------------

\subsubsection*{28. Exemple avec plusieurs zéros}

On peut prendre
\[
h(x)=x(x-1)
\]
sur
\[
[0,1].
\]
Alors
\[
h(0)=0
\qquad\text{et}\qquad
h(1)=0.
\]
La condition
\[
h(0)\leq0\leq h(1)
\]
est vérifiée.

La fonction \(h\) possède plusieurs zéros dans \([0,1]\), à savoir
\[
0
\quad\text{et}\quad
1.
\]

Un autre exemple plus riche est
\[
h(x)=x^2-\frac14
\]
sur \([-1,1]\). Alors
\[
h(-1)>0,
\qquad
h(0)<0,
\qquad
h(1)>0,
\]
et les zéros sont
\[
-\frac12
\quad\text{et}\quad
\frac12.
\]
Mais l'exemple \(x(x-1)\) suffit.

% ------------------------------------------------------------

\subsubsection*{29. Dans ce cas, la dichotomie converge-t-elle forcément vers un zéro ?}

Oui, avec la procédure telle qu'elle a été définie, la dichotomie converge toujours vers un zéro de \(h\), même si \(h\) possède plusieurs zéros.

En effet, la preuve de la partie A n'utilise pas l'unicité du zéro. Elle montre que les suites \((a_n)\) et \((b_n)\) sont adjacentes et ont une limite commune \(\alpha\). Elle montre ensuite, grâce à la continuité de \(h\), que
\[
h(\alpha)=0.
\]
Donc la limite construite est forcément un zéro.

En revanche, si plusieurs zéros existent, la procédure ne permet pas nécessairement de choisir à l'avance lequel elle va approcher. Le zéro obtenu dépend des signes rencontrés aux milieux successifs et donc des choix imposés par la règle de dichotomie.

% ------------------------------------------------------------

\subsubsection*{30. Pourquoi la dichotomie est-elle robuste ?}

La dichotomie est robuste parce qu'elle repose uniquement sur trois idées simples :

\begin{itemize}
    \item la continuité de la fonction ;
    \item un changement de signe ;
    \item le découpage successif d'un intervalle en deux.
\end{itemize}

Elle ne demande pas de résoudre explicitement l'équation
\[
h(x)=0.
\]
Elle ne demande pas non plus de deviner une formule exacte pour la racine.

À chaque étape, on obtient un intervalle qui contient une racine, avec une longueur connue. Ainsi, on contrôle rigoureusement l'erreur.

\begin{resultatbox}
La dichotomie transforme une preuve d'existence en méthode d'approximation certifiée.
\end{resultatbox}

\newpage

% ============================================================
% SYNTHESE
% ============================================================

\section*{Synthèse finale du corrigé}

\subsection*{Méthodes majeures utilisées}

\begin{itemize}
    \item La borne supérieure permet de construire un réel dont l'existence n'est pas donnée explicitement.
    \item Une suite récurrente \(u_{n+1}=f(u_n)\) s'étudie souvent par stabilité d'un intervalle, monotonie et passage à la limite.
    \item Le théorème des accroissements finis et l'inégalité des accroissements finis permettent d'obtenir des inégalités sans développement limité.
    \item Le théorème de Rolle peut servir à montrer qu'une fonction ne prend pas deux fois la même valeur.
    \item Le théorème de la bijection continue permet de justifier proprement l'existence d'une fonction réciproque.
    \item La dichotomie donne une version constructive du théorème des valeurs intermédiaires.
\end{itemize}

\subsection*{Résultats principaux}

\[
\sup\{x\in\R_+ \mid x^2\leq2\}=\sqrt2.
\]

La suite de Héron
\[
u_{n+1}=\frac12\left(u_n+\frac2{u_n}\right),
\qquad u_0=2,
\]
converge vers
\[
\sqrt2.
\]

La fonction
\[
f(x)=\frac{\sqrt{1+x}-1}{x}
\]
prolongée par
\[
f(0)=\frac12
\]
est continue et dérivable sur \(]-1,+\infty[\), strictement décroissante, et réalise une bijection de \(]-1,+\infty[\) sur \(]0,1[\).

Sa réciproque est
\[
f^{-1}(y)=\left(\frac1y-1\right)^2-1.
\]

Pour \(x>0\),
\[
0\leq 1+\frac{x}{2}-\sqrt{1+x}\leq \frac{x^2}{8}.
\]

La procédure de dichotomie construit des suites adjacentes \((a_n)\) et \((b_n)\), avec
\[
b_n-a_n=\frac{b-a}{2^n},
\]
et leur limite commune est un zéro de la fonction continue considérée.

\end{document}
